
% =========================================================== 8. Caso Brasil, os dados
\begin{frame}
  \framesubtitle{Caso Brasil \ \textbullet\ \ quarta, 4 de fevereiro de 2026}
  \frametitle{O parque sobe 43 GW e o CMO se multiplica por 8.}

  \begin{columns}[T,onlytextwidth]
  \begin{column}{0.545\textwidth}
    \vspace{0.5em}
    \centering
    \resizebox{\linewidth}{!}{\begin{tikzpicture}[x=1cm,y=1cm]
  \draw[rulec,line width=.6pt] (0,0) -- (7.15,0);
  \draw[rulec,line width=.6pt] (0,0) -- (0,2.85);
  \draw[rulec,line width=.6pt] (7.15,0) -- (7.15,2.85);
  \draw[rulec,line width=.6pt] (-.07,0.000) -- (0,0.000);
  \node[left,muted,font=\fontsize{5.2}{6.2}\selectfont] at (-.09,0.000) {0};
  \draw[rulec,line width=.6pt] (-.07,0.591) -- (0,0.591);
  \node[left,muted,font=\fontsize{5.2}{6.2}\selectfont] at (-.09,0.591) {25};
  \draw[rulec,line width=.6pt] (-.07,1.182) -- (0,1.182);
  \node[left,muted,font=\fontsize{5.2}{6.2}\selectfont] at (-.09,1.182) {50};
  \draw[rulec,line width=.6pt] (-.07,1.773) -- (0,1.773);
  \node[left,muted,font=\fontsize{5.2}{6.2}\selectfont] at (-.09,1.773) {75};
  \draw[rulec,line width=.6pt] (-.07,2.364) -- (0,2.364);
  \node[left,muted,font=\fontsize{5.2}{6.2}\selectfont] at (-.09,2.364) {100};
  \draw[rulec,line width=.6pt] (7.15,0.000) -- (7.22,0.000);
  \node[right,brick,font=\fontsize{5.2}{6.2}\selectfont] at (7.24,0.000) {0};
  \draw[rulec,line width=.6pt] (7.15,1.300) -- (7.22,1.300);
  \node[right,brick,font=\fontsize{5.2}{6.2}\selectfont] at (7.24,1.300) {2,5k};
  \draw[rulec,line width=.6pt] (7.15,2.600) -- (7.22,2.600);
  \node[right,brick,font=\fontsize{5.2}{6.2}\selectfont] at (7.24,2.600) {5k};
  \fill[brick,opacity=.17] (-0.096,0) rectangle (0.096,0.191);
  \fill[brick,opacity=.17] (0.215,0) rectangle (0.407,0.177);
  \fill[brick,opacity=.17] (0.525,0) rectangle (0.718,0.171);
  \fill[brick,opacity=.17] (0.836,0) rectangle (1.029,0.170);
  \fill[brick,opacity=.17] (1.147,0) rectangle (1.340,0.170);
  \fill[brick,opacity=.17] (1.458,0) rectangle (1.651,0.171);
  \fill[brick,opacity=.17] (1.769,0) rectangle (1.962,0.170);
  \fill[brick,opacity=.17] (2.080,0) rectangle (2.272,0.166);
  \fill[brick,opacity=.17] (2.391,0) rectangle (2.583,0.166);
  \fill[brick,opacity=.17] (2.701,0) rectangle (2.894,0.163);
  \fill[brick,opacity=.17] (3.012,0) rectangle (3.205,0.163);
  \fill[brick,opacity=.17] (3.323,0) rectangle (3.516,0.163);
  \fill[brick,opacity=.17] (3.634,0) rectangle (3.827,0.160);
  \fill[brick,opacity=.17] (3.945,0) rectangle (4.138,0.147);
  \fill[brick,opacity=.17] (4.256,0) rectangle (4.449,0.166);
  \fill[brick,opacity=.17] (4.567,0) rectangle (4.759,0.171);
  \fill[brick,opacity=.17] (4.878,0) rectangle (5.070,0.179);
  \fill[brick,opacity=.17] (5.188,0) rectangle (5.381,0.186);
  \fill[brick,opacity=.17] (5.499,0) rectangle (5.692,0.301);
  \fill[brick,opacity=.17] (5.810,0) rectangle (6.003,1.311);
  \fill[brick,opacity=.17] (6.121,0) rectangle (6.314,2.513);
  \fill[brick,opacity=.17] (6.432,0) rectangle (6.625,1.615);
  \fill[brick,opacity=.17] (6.743,0) rectangle (6.936,0.454);
  \fill[brick,opacity=.17] (7.054,0) rectangle (7.246,0.203);
  \draw[muted,line width=.45pt] (0.000,1.993) -- (0.311,1.912) -- (0.622,1.855) -- (0.933,1.827) -- (1.243,1.823) -- (1.554,1.858) -- (1.865,1.909) -- (2.176,2.036) -- (2.487,2.175) -- (2.798,2.248) -- (3.109,2.307) -- (3.420,2.335) -- (3.730,2.297) -- (4.041,2.352) -- (4.352,2.393) -- (4.663,2.387) -- (4.974,2.388) -- (5.285,2.290) -- (5.596,2.273) -- (5.907,2.348) -- (6.217,2.359) -- (6.528,2.347) -- (6.839,2.277) -- (7.150,2.163);
  \draw[tealx,line width=.65pt] (0.000,0.141) -- (0.311,0.130) -- (0.622,0.131) -- (0.933,0.130) -- (1.243,0.134) -- (1.554,0.143) -- (1.865,0.239) -- (2.176,0.490) -- (2.487,0.674) -- (2.798,0.814) -- (3.109,0.923) -- (3.420,0.977) -- (3.730,0.976) -- (4.041,0.928) -- (4.352,0.830) -- (4.663,0.699) -- (4.974,0.538) -- (5.285,0.298) -- (5.596,0.142) -- (5.907,0.129) -- (6.217,0.146) -- (6.528,0.139) -- (6.839,0.140) -- (7.150,0.133);
  \draw[ink,line width=1.0pt] (0.000,1.862) -- (0.311,1.793) -- (0.622,1.737) -- (0.933,1.709) -- (1.243,1.703) -- (1.554,1.734) -- (1.865,1.689) -- (2.176,1.538) -- (2.487,1.458) -- (2.798,1.351) -- (3.109,1.286) -- (3.420,1.243) -- (3.730,1.185) -- (4.041,1.289) -- (4.352,1.444) -- (4.663,1.596) -- (4.974,1.764) -- (5.285,1.946) -- (5.596,2.111) -- (5.907,2.201) -- (6.217,2.196) -- (6.528,2.193) -- (6.839,2.127) -- (7.150,2.033);
  \draw[ink,line width=.35pt,dashed] (3.730,0) -- (3.730,2.85);
  \filldraw[fill=paperbg,draw=ink,line width=.7pt] (3.730,1.185) circle (1.3pt);
  \node[left,ink,font=\fontsize{5.6}{6.7}\selectfont\bfseries] at (3.650,2.75) {t1};
  \node[left,brick,font=\fontsize{5.2}{6.2}\selectfont] at (3.650,2.52) {50 GW};
  \draw[ink,line width=.35pt,dashed] (5.907,0) -- (5.907,2.85);
  \filldraw[fill=paperbg,draw=ink,line width=.7pt] (5.907,2.201) circle (1.3pt);
  \node[left,ink,font=\fontsize{5.6}{6.7}\selectfont\bfseries] at (5.827,2.75) {t2};
  \node[left,brick,font=\fontsize{5.2}{6.2}\selectfont] at (5.827,2.52) {93 GW};
  \node[below,muted,font=\fontsize{5.2}{6.2}\selectfont] at (0.000,-.04) {0h};
  \node[below,muted,font=\fontsize{5.2}{6.2}\selectfont] at (0.933,-.04) {3h};
  \node[below,muted,font=\fontsize{5.2}{6.2}\selectfont] at (1.865,-.04) {6h};
  \node[below,muted,font=\fontsize{5.2}{6.2}\selectfont] at (2.798,-.04) {9h};
  \node[below,muted,font=\fontsize{5.2}{6.2}\selectfont] at (3.730,-.04) {12h};
  \node[below,muted,font=\fontsize{5.2}{6.2}\selectfont] at (4.663,-.04) {15h};
  \node[below,muted,font=\fontsize{5.2}{6.2}\selectfont] at (5.596,-.04) {18h};
  \node[below,muted,font=\fontsize{5.2}{6.2}\selectfont] at (6.528,-.04) {21h};
\end{tikzpicture}}

    \vspace{0.5em}
    \raggedright
    {\fontsize{5.9}{7.4}\selectfont\color{muted}
     \textcolor{ink}{\rule[0.22ex]{0.9em}{1.1pt}}~demanda líquida\hfill
     \textcolor{muted}{\rule[0.22ex]{0.9em}{0.5pt}}~carga bruta\hfill
     \textcolor{tealx}{\rule[0.22ex]{0.9em}{0.7pt}}~eólica + solar\hfill
     \textcolor{brick!35}{\rule[0ex]{0.6em}{0.6em}}~CMO, eixo direito\par}

    \vspace{0.4em}
    \gloss{t1 é o vale (12h) e t2 é a ponta (19h) da curva do parque controlável.
    A carga bruta quase não se mexe: quem sobe 43 GW é o parque controlável.}
  \end{column}

  \begin{column}{0.425\textwidth}
    \vspace{0.5em}
    {\fontsize{7.3}{9.2}\selectfont
    \setlength{\tabcolsep}{2.8pt}
    \renewcommand{\arraystretch}{1.12}
    \begin{tabular}{l r r r}
      \toprule
       & \textbf{t1 · 12h} & \textbf{t2 · 19h} & \textbf{$\Delta$} \\
      \midrule
      Carga bruta        & 97.168 & 99.352 & \hl{$+$2.185} \\
      Eólica + solar     & 41.281 & 5.464  & \hl{$-$35.816} \\
      Distribuída        & 28.428 & 278    & \hl{$-$28.150} \\
      \addlinespace[0.1em]\midrule
      Hidráulica         & 41.696 & 81.048 & $+$39.352 \\
      Térmica            & 8.453  & 12.064 & $+$3.611 \\
      \textbf{Controlável} & \hg{50.148} & \hg{93.111} & \hg{$+$42.963} \\
      \addlinespace[0.1em]\midrule
      CMO (SE)           & 309    & \hl{2.521} & $\times$8,2 \\
      PLD (SE)           & 350    & 1.200  & $\times$3,4 \\
      \bottomrule
    \end{tabular}}

    \vspace{0.5em}
    \thesis{\fontsize{7.1}{9}\selectfont
      Dois dias depois o mesmo sistema moveu 27 GW e o CMO andou 47 R\$/MWh.
      No dia 4 moveu 43 GW e o CMO andou 2.212. O valor da flexibilidade está
      concentrado em poucas horas do ano.}
  \end{column}
  \end{columns}
\end{frame}

% =========================================================== 9. Caso Brasil, a pilha
\begin{frame}
  \framesubtitle{Caso Brasil \ \textbullet\ \ a pilha e a bateria}
  \frametitle{A rampa da hidráulica é o que separa os dois preços.}

  \begin{columns}[T,onlytextwidth]
  \begin{column}{0.545\textwidth}
    \vspace{0.5em}
    \centering
    \resizebox{\linewidth}{!}{\input{figures/fig-pilha.tikz}}

    \vspace{0.5em}
    \raggedright
    {\fontsize{5.9}{7.4}\selectfont\color{muted}
     \textcolor{muted!60}{\rule[0ex]{0.6em}{0.6em}}~hidro inflexível\hfill
     \textcolor{tealx!45}{\rule[0ex]{0.6em}{0.6em}}~hidro flexível\hfill
     \textcolor{brick!40}{\rule[0ex]{0.6em}{0.6em}}~térmicas por CVU\par
     \textcolor{ink}{\rule[0.22ex]{0.9em}{0.4pt}}~quantidade da hora\hfill
     \textcolor{brick}{\rule[0.22ex]{0.9em}{0.8pt}}~preço do modelo\hfill\mbox{}\par}

    \vspace{0.4em}
    \gloss{As verticais tracejadas são quantidades, o que o parque entregou em
    cada hora; as horizontais são preços. Elas se afastam quando a hidráulica
    flexível deixa de alcançar a ponta. Rampa de 14,3\%/h da flexível, a do dia.}
  \end{column}

  \begin{column}{0.425\textwidth}
    \vspace{0.5em}
    {\fontsize{6.5}{8.4}\selectfont
    \setlength{\tabcolsep}{1.9pt}
    \renewcommand{\arraystretch}{1.12}
    \begin{tabular}{l r r r @{\hspace{0.55em}} r r r}
      \toprule
      {\fontsize{5.6}{6}\selectfont R\$ mi}
        & \multicolumn{3}{c}{\textbf{Sem bat. \ 309\,$\rightarrow$\,446}}
        & \multicolumn{3}{c}{\textbf{Com bat. \ 309\,$\rightarrow$\,347}} \\
      \cmidrule(r){2-4}\cmidrule(l){5-7}
       & Rec. & Custo & Lucro & Rec. & Custo & Lucro \\
      \midrule
      Hidro inflex. & 31,5 & 0,0 & 31,5 & 27,3 & 0,0 & 27,3 \\
      Hidro flexível   & 17,7 & 12,3 & 5,4 & 14,4 & 12,9 & 1,5 \\
      Térmica          & 7,9 & 4,3 & 3,6 & 6,0 & 3,4 & 2,6 \\
      Bateria          & & & & \hg{0,1} & \hg{0,0} & \hg{0,1} \\
      \addlinespace[0.1em]\midrule
      Demanda líq. & \hl{$-$57,0} & & & \hl{$-$47,8} & & \\
      \addlinespace[0.1em]\midrule
      \textbf{Soma}    & \textbf{0,0} & \textbf{16,5} & \textbf{40,5}
                       & \textbf{0,0} & \textbf{16,3} & \textbf{31,5} \\
      \bottomrule
    \end{tabular}}

    \vspace{0.35em}
    \gloss{Com 2.000 MWh de bateria. A receita soma zero nos dois casos.}

    \vspace{0.4em}
    \thesis{\fontsize{7.1}{9}\selectfont
      A demanda paga R\$ 9,2 milhões a menos, e o custo de operação cai só
      R\$ 0,2 milhão: quase tudo é transferência dos inframarginais para o
      consumidor, e a bateria fica com R\$ 0,1 milhão.}
  \end{column}
  \end{columns}
\end{frame}
